% !TEX root = ../main.tex

\section{任务设计与评测方案}

\subsection{任务库设计原则}

课程要求每个网站至少 3 个任务。我们没有只选择简单搜索任务，而是按 easy、medium、hard 三档构建任务库，使任务难度逐步增加：easy 任务侧重页面定位和单点信息抽取等简单任务；medium 任务往往需要跨页面查找、比较和结构化字段补全；hard 任务强调长链推理、跨来源验证、候选集合筛选和最终报告生成。

核心评测包含 CSRankings 两项、Bilibili 三项与广州旅游三项任务，合计八项，分别存于三个主线 YAML 文件。另将 CSRankings 跨方向统计单独存于探索任务目录；由于题目与判分设置需要进一步审查，该任务不进入核心完成率。每项任务均定义任务标识、网站、难度、起始地址、自然语言说明、最大步数、成功标准和禁止动作。

\begin{table}[H]
  \centering
  \caption{8 项核心评测任务与 1 项补充探索任务}
  \label{tab:task-suite}
  \footnotesize
  \setlength{\tabcolsep}{4pt}
  \begin{tabularx}{\textwidth}{p{0.15\textwidth}p{0.10\textwidth}p{0.18\textwidth}Y}
    \toprule
    场景 & 难度 & task id 摘要 & 任务摘要 \\
    \midrule
    CSRankings & easy & NLP professor & 在 2020--2025 年美国 NLP 排名最高学校中寻找论文数最多、排名最靠前的教授，确定性答案为 Graham Neubig。\\
    CSRankings & medium & CN ML scholar & 找 2015--2025 年中国大陆高校 Machine Learning 排名最高学校中排名最高学者，并补全职位、邮箱、引用量、H-index 和最高引论文标题。\\
    CSRankings & 探索 & AI/Sys. faculty & 不计入核心指标；比较美国高校在 8 个 AI/Systems 方向前 30 名中的均衡性，计算中位数、平均值和最差排名，再找覆盖方向最广的 faculty。\\
    Bilibili & easy & public video & 从 B 站主页进入公开视频页，返回标题、链接和可见简介；若无简介需明确说明。\\
    Bilibili & medium & AI top 3 & 打开“人工智能”搜索结果页，读取综合排序前三个 UP 主视频的可见信息，并分析受众与理由。\\
    Bilibili & hard & bangumi report & 浏览至少 5 个热门原创番剧详情，选 3 个外部站核验，最终选择 1 个推荐对象，获取相关视频下载链接并生成 HTML 展示页。\\
    广州旅游 & easy & top attraction & 在携程查找广州景点页，返回排序第 1 景点的名称、评分、地址或区域、开放时间或门票提示和详情页 URL。\\
    广州旅游 & medium & Chimelong & 只使用携程景点页，至少打开 3 个广州长隆相关景点详情，比较评分、地址、开放时间、票务提示并推荐 1 个。\\
    广州旅游 & hard & two-day plan & 从广州南站进出，基于真实网页规划两天一晚行程；查酒店、餐厅、官方景点页面，给出交通、预算、雨天替换和来源 URL。\\
    \bottomrule
  \end{tabularx}
\end{table}

\subsection{难度划分}

难度不单由步数决定，而由“状态空间、证据要求、判分复杂度”共同决定。CSRankings 核心任务侧重教授定位和跨来源档案补全；补充探索任务涉及 8 个方向的学校排名、统计量计算与 faculty 覆盖统计，当前用于检查题目与判分的适用性。Bilibili 与旅游任务更开放，hard 任务不仅需要浏览多个候选页面，还要跨来源核验和生成最终展示材料。

\begin{table}[H]
  \centering
  \caption{Easy/Medium/Hard 分层标准}
  \label{tab:difficulty}
  \small
  \begin{tabularx}{\textwidth}{p{0.14\textwidth}p{0.22\textwidth}Y}
    \toprule
    难度 & 典型能力 & 评测重点 \\
    \midrule
    Easy & 目标明确，少量页面跳转，单点信息抽取 & 元素定位、滚动、基础阅读、答案是否精确。\\
    Medium & 多页面跳转、字段补全、候选比较 & 跨页面记忆、来源记录、结构化输出完整性。\\
    Hard & 长链任务、跨来源验证、复杂计算或报告生成 & 计划分解、证据覆盖、避免循环、开放结果可复核性。\\
    \bottomrule
  \end{tabularx}
\end{table}

\subsection{评测指标与判分口径}

批量运行后，\texttt{eval/metrics.py} 从每个 \texttt{run\_summary.json} 和 step trace 构造 run record，并计算以下指标：
\begin{itemize}
  \item \textbf{Task Success Rate (SR)}：harness 判定成功的 run 数 / 总 run 数，反映主循环是否顺利停止并生成答案；
  \item \textbf{Validated Success Rate}：对确定性任务按答案规则复核后的成功率；
  \item \textbf{Average Steps / Average Time}：平均动作步数与耗时；
  \item \textbf{Action JSON Validity Rate}：解析成功且未使用 fallback 的动作步数占比；
  \item \textbf{Grounding Error Rate}：因 \jsonkey{target\_id} 不存在、不可见或不可用导致的定位错误占比；
  \item \textbf{Correction/Retry Success Rate}：触发重试或恢复机制后最终成功的比例；
  \item \textbf{Failure Type Distribution}：将失败归类为 perception、planning、execution、verification、website 五类。
\end{itemize}

CSRankings 的两项核心任务属于封闭答案任务，可由脚本严格复核；补充探索任务单独记录，不参与核心指标。Bilibili 与广州旅游规划任务具有开放性，页面内容会随时间、推荐系统、地区和登录状态变化，因此先自动检查字段覆盖和来源 URL，再结合 trace 进行人工或 LLM-as-judge 复核。本文明确区分 \textbf{harness success}（流程正常结束）和 \textbf{validator success}（答案通过严格规则）；二者不能混用。

\begin{figure}[H]
  \centering
  \begin{tikzpicture}[
      font=\small,
      node distance=1.1cm,
      box/.style={draw, rounded corners=6pt, align=center, minimum width=2.75cm, minimum height=0.9cm, fill=softBlue},
      green/.style={box, fill=softGreen},
      orange/.style={box, fill=softOrange},
      arrow/.style={-{Latex[length=2.3mm]}, thick}
    ]
    \node[box] (yaml) {tasks/*.yaml};
    \node[green, right=of yaml] (runner) {batch runner\\repeat n};
    \node[orange, right=of runner] (trace) {per-run trace\\step artifacts};
    \node[box, below=of trace] (metrics) {指标统计};
    \node[green, left=of metrics] (validator) {答案/字段复核};
    \node[orange, left=of validator] (report) {汇总/轨迹报告};
    \draw[arrow] (yaml) -- (runner);
    \draw[arrow] (runner) -- (trace);
    \draw[arrow] (trace) -- (metrics);
    \draw[arrow] (trace) -- (validator);
    \draw[arrow] (validator) -- (metrics);
    \draw[arrow] (metrics) -- (report);
  \end{tikzpicture}
  \caption{批量评测与报告生成流程}
  \label{fig:evaluation-pipeline}
\end{figure}

\subsection{第一轮评测：DeepSeek V4 Pro 八项核心任务重复评测}

DeepSeek V4 Pro 对八项核心任务各运行 10 次，共 80 次，记录到 76 次 harness 完成，完成率为 \textbf{95.00\%}。CSRankings 两项核心任务共 20 次均完成；Bilibili 三项任务共 30 次均完成；广州旅游三项任务共 30 次完成 26 次。四次未完成分别为广州 medium 的三次规划失败和广州 hard 的一次感知失败。

原始实验还对补充探索任务运行了 10 次。由于该题的任务定义与判分设置需要进一步核验，这部分记录与八项核心任务分开，不用于推断模型成功率。原九项混合批次的平均时间、平均步数及失败比例也不直接沿用为核心任务指标。

\begin{table}[H]
  \centering
  \caption{DeepSeek V4 Pro 正式评测场景级结果}
  \label{tab:deepseek-scenario}
  \small
  \begin{tabularx}{\textwidth}{p{0.24\textwidth}p{0.14\textwidth}p{0.14\textwidth}p{0.14\textwidth}Y}
    \toprule
    场景 & 运行次数 & 成功次数 & 成功率 & 说明 \\
    \midrule
    CSRankings & 20 & 20 & 100.00\% & 两项核心任务各运行 10 次，均完成。\\
    广州旅游规划 & 30 & 26 & 86.67\% & medium 失败 3 次，hard 失败 1 次。\\
    Bilibili & 30 & 30 & 100.00\% & 三个难度均成功。\\
    合计 & 80 & 76 & 95.00\% & 核心任务未完成 4 次。\\
    \bottomrule
  \end{tabularx}
\end{table}

\begin{table}[H]
  \centering
  \caption{DeepSeek V4 Pro 正式评测任务级结果}
  \label{tab:deepseek-repeat}
  \small
  \setlength{\tabcolsep}{4pt}
  \begin{tabularx}{\textwidth}{p{0.18\textwidth}Yp{0.08\textwidth}p{0.10\textwidth}p{0.10\textwidth}p{0.11\textwidth}p{0.13\textwidth}}
    \toprule
    场景/难度 & 任务 & 次数 & 成功 & 成功率 & 平均步数 & 备注 \\
    \midrule
    CS easy & NLP top professor & 10 & 10 & 100\% & 3.4 & 平均用时 186.19 秒。\\
    CS medium & CN ML scholar profile & 10 & 10 & 100\% & 8.1 & 平均用时 308.29 秒。\\
    GZ easy & Ctrip top attraction & 10 & 10 & 100\% & 5.1 & 平均用时 235.94 秒。\\
    GZ medium & Chimelong sights comparison & 10 & 7 & 70\% & 10.7 & 规划失败 3 次。\\
    GZ hard & two-day real web plan & 10 & 9 & 90\% & 34.1 & 感知失败 1 次，平均用时 1119.67 秒。\\
    Bili easy & public video intro & 10 & 10 & 100\% & 2.2 & 平均用时 235.37 秒。\\
    Bili medium & AI search audience & 10 & 10 & 100\% & 2.6 & 平均用时 287.21 秒。\\
    Bili hard & bangumi HTML recommendation & 10 & 10 & 100\% & 26.25 & 平均用时 708.60 秒。\\
    \bottomrule
  \end{tabularx}
\end{table}

核心任务结果表明，CSRankings 两项信息检索、广州 easy 与 Bilibili 三项任务的完成过程较稳定；广州 medium/hard 仍暴露跨页面导航、信息覆盖与证据管理的问题。长任务还具有更高的操作成本，例如广州 hard 平均 34.1 步、1119.67 秒，Bilibili hard 平均 26.25 步、708.60 秒。完成率应与任务范围、步数、用时和校验方式共同解读。

\subsection{GPT-5.5 八项核心任务视觉基线}

GPT-5.5 以 vision observation mode 对八项核心任务各运行一次，均记录为完成，平均步数为 9.75。两项 CSRankings 核心任务通过确定性答案校验；六项开放任务的字段覆盖均为 1.0，内容真实性仍需结合 trace 复核。该基线每题仅一次，不能等同于重复运行可靠性或严格答案正确率。原始混合批次还包含一项探索任务的单次运行；该运行不进入表中核心指标。由于现有汇总没有提供完整的逐任务用时，不为核心子集复用原混合批次的平均用时。

\begin{table}[H]
  \centering
  \caption{GPT-5.5 八项核心任务基线结果}
  \label{tab:gpt55-baseline}
  \small
  \begin{tabularx}{\textwidth}{p{0.33\textwidth}p{0.12\textwidth}p{0.14\textwidth}p{0.11\textwidth}Y}
    \toprule
    任务 & 难度 & success & 步数 & 复核结论 \\
    \midrule
    CSRankings NLP professor & easy & 成功 & 3 & validator 通过，答案为 Graham Neubig。\\
    CSRankings ML profile & medium & 成功 & 15 & validator 通过，朱军档案字段匹配。\\
    Bilibili public video intro & easy & 成功 & 2 & 字段覆盖 1.0，需结合 trace 复核具体视频页。\\
    Bilibili AI search top 3 audience & medium & 成功 & 1 & 字段覆盖 1.0，从搜索页可见信息完成分析。\\
    Bilibili bangumi HTML & hard & 成功 & 20 & 字段覆盖 1.0，完成候选浏览、外部核验和 HTML 输出。\\
    Ctrip top attraction & easy & 成功 & 4 & 字段覆盖 1.0，返回景点、评分、地址与详情页。\\
    Ctrip Chimelong compare & medium & 成功 & 6 & 字段覆盖 1.0，完成多景点详情页对比。\\
    Guangzhou two-day plan & hard & 成功 & 27 & 字段覆盖 1.0，完成酒店、餐饮、景点与路线规划。\\
    \bottomrule
  \end{tabularx}
\end{table}

\subsection{十模型横向评测}

我们随后选取 10 个多样化模型进行横向评测，覆盖 OpenAI、Google Gemini、Anthropic Claude、Qwen、DeepSeek、Kimi、GLM 和豆包系列。视觉模型为 GPT-4o mini、GPT-5 mini、Gemini 2.5 Flash、Claude Haiku 4.5、Qwen3-VL-32B 与豆包 Vision Pro；文本模型为 DeepSeek-V4-Pro、Qwen3.7-Max、Kimi-K2.6 与 GLM-5.2。完整 API 名称见表~\ref{tab:model-valid-results}。前 6 个按 vision mode 运行，后 4 个按 text mode 运行。历史批次同时包含八项核心任务和一项补充探索任务，与 10 个模型构成 90 个计划组合。核心评测的设计范围是 80 个模型--任务组合，探索部分另有 10 个；计划规模不等于有效样本数。

我们在此次评测中也遇到了诸如中转站超时、内容安全拦截、上下文超限等服务层问题，为了避免误计为模型能力失败，我们把每个组合标记为 \textbf{valid}、\textbf{invalid} 或 \textbf{incomplete}。表~\ref{tab:model-valid-results} 保留历史混合批次的汇总，用于追溯原始实验覆盖：CS 列同时包含两项核心任务和一项探索任务，开放列包含六项核心任务，数对按“有效样本数/通过或完成数”排列。汇总合计为 66 个有效、22 个无效和 2 个未完成。该表未逐项标明部分模型的探索任务状态，不能据此精确重建八任务核心子集，也不能将 CS 列直接用作核心任务排行榜。

\begingroup
\small
\renewcommand{\arraystretch}{1.15}
\setlength{\tabcolsep}{4pt}
\begin{longtable}{p{0.15\textwidth}p{0.07\textwidth}p{0.10\textwidth}p{0.11\textwidth}p{0.065\textwidth}p{0.07\textwidth}p{0.265\textwidth}}
  \caption{十模型历史批次记录：核心任务与补充探索任务的混合汇总}\label{tab:model-valid-results}\\
  \toprule
  模型 & 模式 & \makecell[c]{CS\\严格成功} & \makecell[c]{开放\\完成} & 无效 & \makecell[c]{未\\完成} & 有效样本中的主要观察 \\
  \midrule
  \endfirsthead
  \toprule
  模型 & 模式 & \makecell[c]{CS\\严格成功} & \makecell[c]{开放\\完成} & 无效 & \makecell[c]{未\\完成} & 有效样本中的主要观察 \\
  \midrule
  \endhead
  \bottomrule
  \endfoot
  gpt-4o-mini & vision & 3/0 & 6/2 & 0 & 0 & 探索任务和部分开放长任务出现动作循环。\\
  gpt-5-mini & vision & 3/2 & 6/4 & 0 & 0 & 核心 easy/medium 较稳；探索任务与广州 hard 需进一步检查。\\
  gemini-2.5-flash & vision & 3/1 & 6/3 & 0 & 0 & 可完成部分阅读任务，但在广州 medium/hard 与 B 站 hard 中出现证据收集循环。\\
  claude-haiku-4-5 & vision & 3/1 & 0/0 & 6 & 0 & 六项开放任务输入超限；CS 记录含探索任务，不能直接按核心范围比较。\\
  qwen3-vl-32b & vision & 3/2 & 6/4 & 0 & 0 & CS easy/medium 与 B 站任务较好，广州 medium/hard 的真实网页导航仍不稳定。\\
  doubao-vision-pro & vision & 0/0 & 0/0 & 9 & 0 & 图文总 token 超限，当前配置下无可比较有效样本。\\
  DeepSeek-V4-Pro & text & 3/2 & 6/5 & 0 & 0 & 纯文本模式表现稳健，说明 DOM/A11y/page text 对信息型任务有较强支撑。\\
  qwen3.7-max & text & 2/2 & 5/5 & 2 & 0 & 有效样本均完成，但广州 hard 出现内容安全拦截，需要补跑。\\
  kimi-k2.6 & text & 2/2 & 3/2 & 2 & 2 & 有效样本尚可，但服务超时和长任务未完成影响正式统计。\\
  GLM-5.2 & text & 2/2 & 4/3 & 3 & 0 & easy 与 B 站 medium 较好；广州 medium 超时、广州 hard 被内容拦截。\\
\end{longtable}
\endgroup

历史实验提供以下定性观察。第一，流程完成与答案正确需要分别记录；确定性核心任务采用规则校验，开放任务还需结合字段与来源证据复核，补充探索任务不用于核心成功率推断。第二，文本观察中的 DOM、A11y 和页面文本可以支持信息型任务，但当前视觉与文本模式的输入预算不同，不能据此推断视觉输入的因果效果。第三，开放长任务中的重复导航、过早回答和来源遗漏表明，规划与证据管理仍需改进。上述观察用于系统诊断，不构成八项核心任务上的完整模型排名。

\subsection{十模型评测中的无效样本与风险}

在报告有效样本结论之后，还必须单独说明不可归因为模型能力的技术问题。第一类是 \textbf{API read timeout}，如 \texttt{kimi-k2.6} 的 Bilibili medium、\texttt{GLM-5.2} 的广州 medium：请求已发出，但中转站或上游模型在超时时间内没有返回。第二类是 \textbf{HTTP 400 内容安全拦截}，如 \texttt{qwen3.7-max} 和 \texttt{GLM-5.2} 在广州 hard 中被上游安全策略拦截；真实网页文本中下载链接、支付、登录、预约等敏感词可能共同触发代理侧策略，不能直接解释为规划能力差。第三类是 \textbf{上下文或图文 token 超限}：\texttt{claude-haiku-4-5} 的 6 个开放任务、\texttt{doubao-1-5-vision-pro-32k} 的多个任务因此无效。

造成 token 超限的直接原因是当前 vision mode 的 observation 较重：每一步会同时发送任务、动作历史、候选元素、页面文本、Accessibility Tree、raw screenshot 与 marked screenshot。以 Bilibili easy 为例，Claude 的第一步 prompt 文本约 55 KB、Accessibility Tree 约 36 KB、候选元素约 25 KB，两张截图接近 1.9 MB；经 base64 和代理侧估算后约为 428k tokens，超过上下文限制。我们当前的 observation space 内容非常丰富，一方面这可以为强上下文模型提供更加丰富的信息，但是另一方面这也导致了短上下文模型输入超限而造成任务失败。后续如果想评测此类模型，可能需要专门针对其简化observation space 设计。

